% Módulo 05 — Avaliação empírica e testes numéricos (R762+; números integrais)
\chapter{Avaliação empírica e testes numéricos}\label{cap:resultados}
Esta seção apresenta os resultados históricos da bancada e das ablações, distinguidos pelos conjuntos de instâncias descritos na Seção~\ref{cap:metodo}.
\section{Fuzz em ponto flutuante}
O relatório interno registra testes aleatórios de 20\,000 vetores, sem violações acima das tolerâncias examinadas: $10^{-9}$ para a soma e $10^{-12}$ para o deslocamento. O maior desvio registrado da soma foi aproximadamente $2{,}2\times10^{-16}$ e o do deslocamento foi $0{,}0$. A preservação de argmax, os limites, o caso de candidato único e o tratamento do conjunto vazio também foram examinados. A ausência de violações nessa amostra não demonstra correção para todas as entradas. Testes por entradas aleatórias têm antecedentes na avaliação de utilitários Unix \cite{miller1990}.
\section{Bancada principal e variação do parâmetro do oráculo}
Bancada principal (200 decisões): atenção $0{,}042\pm0{,}011$ (acerto $0{,}625\pm0{,}067$) contra $0{,}19$--$0{,}22$; latência $2{,}17$\,ms; soma softmax 1 até $1{,}1\times10^{-16}$ Os valores da bancada principal constam de \nolinkurl{resultado_bancada.json}, reproduzido no Apêndice~\ref{ap:json}.
A Tabela~\ref{tab:bench} apresenta as médias das ablações para cada valor do parâmetro.
\begin{table}[htbp]
\fontsize{10}{12}\selectfont\SingleSpacing
\centering
\caption{Ablações: arrependimento por parâmetro do oráculo e diferença em relação à atenção (200 decisões por configuração).}
\label{tab:bench}
\begin{tabular}{lccc}
\toprule
$q_w$ & Atenção & Pesos iguais (dif.) & Sem confiança (dif.) \\
\midrule
0,3 & $0{,}0146$ & $0{,}0319$ (+0,017) & $0{,}1072$ (+0,093) \\
0,6 & $0{,}0292$ & $0{,}0638$ (+0,035) & $0{,}2144$ (+0,185) \\
0,9 & $0{,}0438$ & $0{,}0957$ (+0,052) & $0{,}3216$ (+0,278) \\
\bottomrule
\end{tabular}
\fonte{Elaboração própria (2026), a partir dos artefatos e das fontes identificados no texto.}
\end{table}
Outras condições de ablação ($q_w=0{,}6$): confiança embaralhada $0{,}2243$ (+0,195), anticorrelacionada $0{,}3604$ (+0,331), ruidosa $0{,}1515$ (+0,122); vitória/empate da atenção $\geq 0{,}92$ em todos os pareamentos. Os registros internos relatam reexecução determinística; uma reprodução independente permanece pendente. A Tabela~\ref{tab:pareado} reúne as diferenças pareadas e as taxas de vitória ou empate.
\begin{table}[htbp]
\fontsize{10}{12}\selectfont\SingleSpacing
\centering
\caption{Diferenças pareadas por decisão vs.\ atenção ($q_w=0{,}6$) e taxa de vitória/empate.}
\label{tab:pareado}
\begin{tabular}{lcc}
\toprule
Estratégia & Diferença $\pm$ margem exploratória & Atenção vence/empata \\
\midrule
Média igualitária & $+0{,}035\pm0{,}011$ & $0{,}99$ \\
Sem confiança & $+0{,}185\pm0{,}024$ & $0{,}98$ \\
Confiança ruidosa & $+0{,}122\pm0{,}022$ & $0{,}93$ \\
Confiança embaralhada & $+0{,}195\pm0{,}025$ & $0{,}92$ \\
Confiança anticorrelacionada & $+0{,}331\pm0{,}019$ & $0{,}97$ \\
\bottomrule
\end{tabular}
\fonte{Elaboração própria (2026), a partir dos artefatos e das fontes identificados no texto.}
\end{table}
\begin{table}[htbp]
\fontsize{10}{12}\selectfont\SingleSpacing
\centering
\caption{Arrependimento por estratégia nos três valores do parâmetro do oráculo (200 decisões cada).}
\label{tab:mundos}
\begin{tabular}{lccc}
\toprule
Estratégia & $q_w=0{,}3$ & $q_w=0{,}6$ & $q_w=0{,}9$ \\
\midrule
Atenção & $0{,}0146$ & $0{,}0292$ & $0{,}0438$ \\
Média igualitária & $0{,}0319$ & $0{,}0638$ & $0{,}0957$ \\
Sem confiança & $0{,}1072$ & $0{,}2144$ & $0{,}3216$ \\
Confiança ruidosa & $0{,}0757$ & $0{,}1515$ & $0{,}2272$ \\
Confiança embaralhada & $0{,}1121$ & $0{,}2243$ & $0{,}3364$ \\
Confiança anticorrelacionada & $0{,}1802$ & $0{,}3604$ & $0{,}5406$ \\
\bottomrule
\end{tabular}
\fonte{Elaboração própria (2026), a partir dos artefatos e das fontes identificados no texto.}
\end{table}
O ranking das estratégias permanece estável na Tabela~\ref{tab:mundos} porque, entre elegíveis com cobertura unitária, o arrependimento é proporcional a $q_w$. As configurações não representam avaliações independentes de oráculos distintos. As diferenças pareadas e suas margens descrevem as instâncias das ablações, sem correção para agrupamento.
\begin{table}[htbp]
\fontsize{10}{12}\selectfont\SingleSpacing
\centering
\caption{Arrependimento por semente e estratégia (40 decisões cada; acerto da atenção entre parênteses).}
\label{tab:seeds}
\begin{tabular}{lcccc}
\toprule
Semente & Atenção & Aleatório & Menor carga & Maior cobertura \\
\midrule
11 & $0{,}0752$ (0,45) & $0{,}1889$ & $0{,}3157$ & $0{,}2319$ \\
22 & $0{,}0172$ (0,78) & $0{,}2463$ & $0{,}2276$ & $0{,}2310$ \\
33 & $0{,}0337$ (0,55) & $0{,}1719$ & $0{,}2447$ & $0{,}2033$ \\
44 & $0{,}0716$ (0,60) & $0{,}1992$ & $0{,}2038$ & $0{,}2309$ \\
55 & $0{,}0129$ (0,75) & $0{,}1350$ & $0{,}0918$ & $0{,}1618$ \\
\bottomrule
\end{tabular}
\fonte{Elaboração própria (2026), a partir dos artefatos e das fontes identificados no texto.}
\end{table}
A atenção obteve o menor arrependimento nas 5 sementes avaliadas (Tabela~\ref{tab:seeds}).
\section{Estudo de caso completo: Ana contra Cid}
Na tarefa-guia, as cabeças mediram: semântica Ana $0{,}8417$ contra Cid $0{,}7703$; capacidade $1{,}0$ para ambas (o gate já havia decidido); confiança $0{,}9$ contra $0{,}7$; carga livre $0{,}9$ contra $0{,}6$. Utilidades: Ana $0{,}30\cdot0{,}8417+0{,}35\cdot1{,}0+0{,}25\cdot0{,}9+0{,}10\cdot0{,}9=0{,}9175$; Cid $0{,}30\cdot0{,}7703+0{,}35\cdot1{,}0+0{,}25\cdot0{,}7+0{,}10\cdot0{,}6=0{,}8161$. Pesos finais $0{,}5253$ contra $0{,}4747$. Note-se a capacidade constante: com cobertura total dos elegíveis, a decisão correu nas outras três cabeças --- exatamente o fenômeno quantificado nas correlações.
\section{Correlações de Pearson (200 decisões)}
Arrependimento $\times$ elegíveis $+0{,}08$; arrependimento $\times$ dispersão $-0{,}16$; acerto $\times$ elegíveis $-0{,}19$; acerto $\times$ dispersão $+0{,}13$; confiança $\times$ $q$ $+0{,}85$ (checagem de desenho: proxy forte, fonte da circularidade parcial); cabeças quase descorrelacionadas (semântica $\times$ confiança $+0{,}05$); capacidade constante após o filtro de elegibilidade (o gate decide cobertura; a escolha corre nas outras cabeças). A transcrição histórica preservada registra 63 testes aprovados em nove suítes. A revisão editorial não reexecutou essa bateria. O diagnóstico preliminar da revisão apresentou 18 verificações aprovadas e duas advertências, sem falhas.
